\section{Challenges and Research Frontiers}
\label{sec:frontiers}

Agentic visual generation is developing along three broad paths: expanding the scope of visual agency, increasing the intelligence available to visual agents, and establishing reliable evaluation. Scope expansion concerns the artifacts, project durations, tools, environments, application consequences, and related operating conditions that a system can handle. Intelligence improvement concerns foundation models, system architectures, reasoning, process training, continual learning, collaboration, deployment, and related capabilities. Evaluation concerns the quality, behavior, cost, and consequences of the complete creation process.

\subsection{How to Expand the Scope of Visual Agency}

\subsubsection{Multimodal Creative Workflows}

A multimodal creative workflow produces several related artifacts during one project. Text, images, audio, video, source code, layouts, data, interactive states, and related representations can all contribute to the final work. Each artifact has its own representation and can influence the content of another artifact. Existing systems often handle each modality through a separate operation with a limited handoff between operations. Multimodal projects are becoming common in storytelling, presentations, education, advertising, scientific communication, interactive media, and other creative settings. Users expect the project to remain coherent when one part changes. A system that produces several outputs without preserving their relations creates additional coordination work for the user. Future systems can represent a project as a set of linked artifacts with explicit dependencies and update rules. A change can be propagated to the artifacts that depend on it, while unrelated artifacts remain stable. The system can maintain shared identities, references, constraints, and versions across modalities. It can also expose which artifacts will change before executing a large update.


\subsubsection{Long-Horizon Creative Projects}

A long-horizon project contains many decisions distributed across turns, scenes, files, production stages, and other production units. The project develops over time, and later actions depend on earlier requirements, intermediate results, user choices, and unresolved issues. Long projects accumulate identity drift, layout drift, inconsistent references, obsolete assumptions, repeated repairs, and similar forms of drift. A system may preserve a conversation while losing the exact artifact version that a decision concerned. It may also continue a plan after the user has changed a central requirement. Long-form video, multi-page documents, interactive stories, product design, complex image editing, and related projects all require continuity across many operations. The user needs the project to remain understandable and editable after a long sequence of changes. Short-task success gives limited information about this ability. Systems can maintain project-level records that connect requirements, artifacts, decisions, and unresolved issues. They can separate stable commitments from provisional choices and update only the dependencies affected by a change. They can preserve checkpoints, compare branches, and present the consequences of a proposed revision before executing it.


\subsubsection{Open Tool Environments}

An open tool environment contains changing models, applications, code runtimes, browsers, simulators, data services, external interfaces, and related services. The available operations, costs, permissions, output formats, and other interface conditions can vary across tasks and over time. Many visual agents rely on a fixed tool library with known interfaces. They may fail when a tool is unavailable, when an output has an unexpected format, when a tool version changes its behavior, or under related environmental changes. They may also select a tool without accounting for execution cost, side effects, or the representation required by a later operation. Visual projects increasingly combine commercial models, open models, desktop applications, code environments, domain-specific software, and related software systems. A fixed toolchain limits the range of tasks that the system can support. An open environment also introduces practical risks that are absent from a closed demonstration. An agent can maintain capability descriptions for tools, learn interface conditions from documentation and execution feedback, and test a tool on a small case before using it in a larger project. It can construct higher-level operations from compatible primitives and record version-specific behavior. Permission boundaries and side-effect descriptions can become part of tool selection.


\subsubsection{Embodied and High-Consequence Domains}

In an embodied or high-consequence domain, a generated artifact supports an action, a decision, a physical process, or a regulated activity. The artifact is evaluated through domain constraints and consequences in addition to visual appearance. Systems often inherit general-purpose evaluation criteria in settings that require specialized evidence. Human responsibility and intervention points may also remain unclear. Visual agents are being applied to navigation, physical simulation, rehabilitation, engineering, education, scientific analysis, public communication, and other high-consequence settings. Errors in these settings can cause physical, financial, medical, or social harm. The acceptable level of autonomy depends on the consequences of an error. Systems can encode domain constraints at the specification and execution levels, monitor evidence during operation, and assign different approval requirements to different actions. They can expose uncertainty and preserve a complete record of the assumptions, references, and checks behind a result. Deployment can use risk-adjusted autonomy with stricter limits for irreversible or high-consequence actions.


\subsection{How to Increase the Intelligence of Visual Agents}

\subsubsection{Generative and Multimodal Foundation Models}

Foundation models provide the perceptual and generative capabilities used by an agent. They determine how well the system can understand an input, produce a visual artifact, edit a region, preserve an identity, and represent information across modalities. A model can produce high-quality individual samples while failing on spatial relations, text, long sequences, precise editing, rare knowledge, physical events, and other task requirements. A model that understands an image may still lack an interface for making an exact change to the source representation. A stronger generator can improve a single operation while leaving the wider control process unchanged. Every higher-level agent decision depends on the quality and controllability of the underlying operations. Weak generation creates repeated repairs and increases the cost of planning. Weak understanding causes the agent to construct an incorrect task representation before any tool is invoked. Foundation models can be improved through richer multimodal pretraining, structured data, synthetic supervision, preference alignment, targeted post-training, unified understanding-generation objectives, and related objectives. Training can include spatial, temporal, structural, factual, physical, and other task constraints. Evaluation should measure which agent decisions improve when the foundation model changes under the same controller and budget.


\subsubsection{Agent Architecture and Control}

Agent architecture determines how task information, visual state, plans, tools, observations, and actions are organized. It includes external orchestration, native visual agents, hierarchical workflows, role-specialized agents, executable programs, combinations of these forms, and related arrangements. A collection of components can still behave as a fixed pipeline. Information can be lost at an interface, a specialist can optimize a local objective that conflicts with the project, and a coordinator can repeat the same action after different failures. The division between planning, execution, observation, and correction may also remain unclear. More complex visual tasks require decisions at several levels: project organization, operation selection, local editing, verification, delivery, and other task decisions. The architecture must support these decisions without creating excessive communication, latency, or error propagation. Future architectures can make state transitions explicit, define interfaces around representations and evidence, and allocate decisions according to the structure of a task. They can combine fast operations for routine cases with deeper search for uncertain cases. They can also allow roles to be created, merged, or suspended according to current task demands.


\subsubsection{Visual Reasoning, World Knowledge, and Causal Understanding}

Visual reasoning connects language, visual evidence, structured relations, external knowledge, possible actions, and related information. Causal understanding adds a prediction of how an artifact or environment will change after an intervention. Systems often produce plausible appearances with incorrect relations, facts, dimensions, event order, or physical consequences. AVG is used for tasks that require spatial precision, factual grounding, procedural structure, and physical plausibility. These requirements cannot be represented by appearance quality alone. An agent must understand which properties are essential and how an action will affect them. Models can combine language and visual reasoning with layouts, scene graphs, executable programs, geometric representations, simulators, retrieval, explicit uncertainty, and related representations. Training can include counterfactual examples that vary an action while holding the initial state fixed. A visual result can then be evaluated for both appearance and its relation to the intended process.


\subsubsection{Agentic Reinforcement Learning and Process Training}

Process training teaches an agent how to make a sequence of decisions. It can supervise tool calls, prompt construction, intermediate judgments, action ordering, corrections, and stopping. Reinforcement learning connects these decisions to outcomes observed later in the trajectory. Final-artifact rewards provide weak information about intermediate decisions. A high score may reward an accidental sample, while a useful repair may receive little credit because its contribution appears several steps later. Learned evaluators can also reward shortcuts that improve a metric while damaging user intent or already-correct regions. Visual creation contains delayed effects, alternative paths, irreversible operations, and competing objectives. Process training can teach the agent to value evidence, preservation, cost, and recovery together with final quality. It can also expose failure patterns that are rare in ordinary demonstration data. Training can combine supervised trajectories, preference comparisons, process rewards, outcome rewards, failure injection, visual feedback, environment feedback, counterfactual branches, and other supervision signals. Agentic RL can learn when to search, inspect, revise, stop, or ask for help. Reward design should account for constraint satisfaction, collateral change, resource use, human intervention, and related objectives.


\subsubsection{Memory, Continual Learning, and Self-Evolution}

Memory makes information from earlier operations available to later decisions. Continual learning changes how the agent behaves across tasks. Self-evolution creates or updates reusable strategies, skills, tools, routing rules, data policies, model parameters, and related learned behaviors. A stored conversation may not identify the artifact version, region, dependency, or evidence associated with a decision. A retrieved example may resemble the current task in language while differing in visual conditions. A self-evolving system may improve a benchmark through extra retrieval or sampling without changing its policy, or it may preserve an error and repeat it across future tasks. Long-term visual work requires continuity, while repeated projects create opportunities for experience reuse. The cost of rediscovering tool paths and repair strategies grows with task complexity. Reliable learning can reduce this cost and support personalization, new tools, and domain transfer. Memory systems can preserve multimodal artifacts, source representations, observations, outcomes, uncertainty, provenance, and related records. Experience can be abstracted into cases, strategies, executable skills, and other reusable forms with explicit preconditions. Updates can be versioned, evaluated on held-out tasks, checked for forgetting and contamination, and reversed after a regression.


\subsubsection{Multi-Agent Collaboration and Human Co-Creation}

Multi-agent collaboration divides interpretation, generation, critique, retrieval, editing, delivery, and related activities across several roles. Human co-creation adds user goals, preferences, contextual knowledge, evaluation, and authorship to the same process. Agents can duplicate work, disagree about the current state, optimize incompatible objectives, and consume resources through communication. Humans can become responsible for checking every low-level operation or lose visibility into important decisions. Creative systems also struggle to balance exploration with coherence and user control. Complex visual creation contains specialized knowledge and open-ended choices. Collaboration can broaden exploration and combine expertise, while human participation can resolve ambiguity and express values that are difficult to formalize. The benefit depends on the quality of coordination and the distribution of initiative. Systems can use shared state, explicit communication protocols, role-specific evidence, conflict resolution, adaptive initiative, user-selectable levels of intervention, and related coordination mechanisms. Agents can propose alternatives, explain trade-offs, and defer high-risk decisions. Evaluation can measure creative diversity, task quality, coordination cost, user workload, authorship, satisfaction with the collaborative process, and related outcomes.


\subsubsection{Efficiency and Scalable Deployment}

Efficient deployment concerns the quality, latency, computational cost, energy, and human time required to operate an agentic visual system. The cost includes model calls, tool calls, retrieval, simulation, communication, candidate generation, review, and related resources. Fixed retry counts treat easy and difficult tasks similarly. More sampling can improve a result while making the workflow impractical. Multi-agent communication and visual verification can dominate the cost of the generator itself. A method that works in a small demonstration may become unusable for long videos, 3D scenes, many users, and other demanding artifacts. A useful agent must allocate effort according to uncertainty and consequence. It should spend more computation on unresolved or high-risk decisions and less on routine operations. Deployment constraints also affect which architectural and learning choices are viable. Systems can estimate the expected value of another action, select verification depth, use early stopping, reuse cached results, and choose between fast and slow paths. Serving infrastructure can optimize multiple stages jointly, while evaluation reports quality--cost and quality--latency curves. Energy, human workload, and related deployment burdens should be included with model and tool expenditure.


\subsection{How to Establish Reliable Evaluation}

\subsubsection{Beyond Final-Artifact Quality}

Final-artifact evaluation measures the appearance, semantic alignment, or task quality of the output. Agent evaluation also measures how the system interpreted the task, chose actions, used evidence, handled resources, and involved people. A final score cannot identify whether improvement came from the foundation model, the controller, additional samples, a human correction, or a lucky branch. It also gives limited information about failed attempts, unused feedback, repeated actions, and hidden intervention. Agentic systems make claims about planning, reflection, memory, recovery, adaptation, and autonomy. These claims concern process behavior. Evaluation that records only the last artifact leaves the central evidence unobserved. Reports can separate artifact quality, constraint satisfaction, state retention, action selection, diagnosis, recovery, stopping, resource use, human intervention, and related process measures. Baselines can hold the generator, tool set, sampling budget, human access, and other operating conditions fixed while changing the controller or learning method.


\subsubsection{Process, Causal, and Counterfactual Evaluation}

Process evaluation records the sequence of states, observations, actions, results, and decisions. Causal evaluation asks whether an action produced the observed improvement. Counterfactual evaluation changes an input, observation, or action and tests the resulting behavior. Many trajectories omit the observation used at a decision point, the version of a tool, the reason for a retry, or the state before a correction. This makes it difficult to determine whether feedback changed the action or whether the system followed a fixed path. The defining property of an agentic loop lies in the connection between evidence and action. A benchmark must expose this connection to distinguish adaptive control from repeated execution. Benchmarks can inject localized defects, unavailable tools, changed references, invalid programs, conflicting constraints, misleading observations, and related perturbations. They can compare the selected action under alternative evidence and measure diagnosis accuracy, repair scope, preserved constraints, recovery time, stopping behavior, and related responses.


\subsubsection{Long-Horizon and Cross-Task Evaluation}

Long-horizon evaluation tests a system over many related decisions. Cross-task evaluation tests whether information or updates from earlier tasks affect later tasks with new prompts, artifacts, tools, users, domains, and related task conditions. Short episodes conceal state drift, memory errors, repeated failures, and stopping problems. A self-improvement claim can also rely on the same task distribution used to create the update, which gives limited evidence for transfer. AVG systems are intended for projects and repeated use. Their value depends on continuity and learning across time. These properties cannot be inferred from a collection of independent single-step outputs. Evaluation sets can separate formation tasks, development tasks, and held-out transfer tasks. They can introduce changed requirements, new tools, new visual domains, delayed references, and related shifts. They can test forgetting, contamination, update persistence, and the relation between a stored experience and a later decision.


\subsubsection{Multidimensional Visual Correctness}

Visual correctness includes several properties that may conflict. These properties include appearance, semantic alignment, spatial relations, geometry, topology, text, data, identity, time, physics, executable behavior, edit preservation, human acceptance, and related requirements. A visually attractive result can contain an incorrect value, missing object, broken relation, spelling error, invalid dimension, identity change, event-order error, or impossible physical consequence. A single score compresses these failures and hides the requirement that caused them. Different visual artifacts expose different forms of correctness. An image, a CAD model, a chart, a page, and an interactive interface cannot be judged through the same evidence source. Agent decisions also depend on knowing which requirement failed. Evaluation can represent requirements as typed claims and associate each claim with an appropriate check. It can combine learned judgments with OCR, parsing, data validation, geometric checks, source execution, simulation, browser interaction, human review, and related evidence sources. Reports can preserve per-claim outcomes and confidence alongside aggregate results.


\subsubsection{Human, Cost, Safety, and Provenance Evaluation}

Human, cost, safety, and provenance evaluation measures the conditions under which a visual agent can be trusted and used. It includes user understanding, intervention timing, workload, resource expenditure, untrusted inputs, tool side effects, privacy, authorship, the history of an artifact, and related factors. User studies often report preference without measuring control or workload. System reports may omit human corrections, hidden API calls, or the cost of generating rejected candidates. Provenance can be lost after editing, delegation, format conversion, or rollback. Visual and audio inputs can also carry instructions that the user did not intend to authorize. Visual agents can affect external files, public content, medical guidance, industrial designs, collaborative projects, and other external settings. Trust depends on a calibrated relationship between system ability, human oversight, and consequence. A quality score alone cannot establish that relationship. Evaluation can report model, tool, retrieval, simulation, communication, candidate, human, and related costs. It can test user understanding, intervention timing, trust calibration, prompt injection, permission boundaries, reversible actions, source attribution, privacy, and related safeguards. Provenance records can connect sources, transformations, responsible actors, approvals, and related events across the full trajectory.


\subsection{More Questions About Agentic Visual Generation}

\subsubsection{What Does Intelligence Mean in Visual Creation?}

Intelligence in visual creation includes understanding, representation, generation, action selection, adaptation, exploration, collaboration, uncertainty management, and responsibility. The relative value of these capacities depends on the task. A system can be strong in visual appearance and weak in structural reasoning. It can be strong in local editing and weak in project continuity. It can be strong in exploration and weak in user control. A single ranking therefore hides important differences between systems. AVG serves personal artwork, scientific figures, engineering designs, educational media, and high-consequence applications. Each setting values a different balance of creativity, accuracy, speed, autonomy, and oversight. Evaluation and system design can describe intelligence as a profile of capabilities under explicit conditions. Reports can state which capacity improved, which constraints were preserved, what resources were used, and how the change affected the user and artifact.


\subsubsection{Authorship, Responsibility, and Human Control}

Agentic creation distributes creative decisions across users, Agents, foundation models, tools, data sources, and evaluators. Authorship concerns contribution and creative direction. Responsibility concerns who can approve, explain, change, or reverse a consequential decision. Users may provide an intention without seeing the intermediate choices that shape the result. Agents may make decisions that affect external files or public content. Multiple models and tools may contribute to one artifact, while their contributions and permissions remain difficult to reconstruct. Creative work involves preferences, cultural context, legal conditions, and personal identity. Human participation can shape the purpose, selection, and interpretation of an artifact. A usable system must support meaningful control at the points where those values matter. Systems can expose editable intermediate states, approval checkpoints, contribution records, provenance, permission boundaries, and reversible actions. They can adapt the level of initiative to the task risk and allow users to inspect assumptions before an irreversible operation.


\subsubsection{From Individual Models to Creative Ecosystems}

A creative ecosystem combines foundation models, Agent policies, tools, data, memory, evaluation services, interfaces, human workflows, and deployment infrastructure. The complete system determines what a user can create and how reliably the result can be revised and delivered. A system can depend on a particular model version, external API, evaluator, data source, or application interface. Changes in one layer can alter the behavior of the complete workflow. Communication overhead, licensing, dependency failures, and attribution become difficult to manage as the ecosystem grows. Visual creation is moving toward reusable workflows and shared services. Progress in one model or tool can have value across many applications when interfaces, provenance, and evaluation are stable. Ecosystem-level design also determines whether research results can be reproduced outside a single demonstration environment. Future systems can specify model and tool contracts, version dependencies, data and provenance records, evaluator behavior, permission boundaries, and resource budgets. They can support replacement of one component while preserving the meaning of the artifact and the trajectory record. The Frontier agenda therefore concerns the changing unit of visual intelligence: a single result, a multimodal workflow, a sustained creation process, a system that learns from its history, and a responsible participant in human creative work. Each stage requires evidence about what the system understood, what it changed, what it learned, what it cost, and how people can control its consequences.


